\section{Forelesning 1: \emph{Introduksjon -- Hva er strategi?}}
\label{sec:lec01}
\settermlabel{terms:lec01}

\subsection*{Oversikt}
Den første forelesningen etablerer det begrepsmessige fundamentet for hele kurset:
hva strategi er, hva strategi \emph{ikke} er, og hvilke antagelser enhver strategi
hviler på. Forelesningen introduserer også kursets gjennomgående linse -- at
IT-strategi handler om bevisste valg om \emph{hvordan} digital teknologi skal brukes
til å posisjonere virksomheten, ikke bare om å digitalisere eksisterende drift.

\begin{figure}[H]
\centering
\includegraphics[width=0.85\linewidth]{lec01-makro-meso-mikro-matrix.png}
\caption{Kursets oversiktsmatrise: hvert tema plasseres på aksene
\emph{makro/meso/mikro} mot \emph{teknologi/forretning}. Brukes til å
posisjonere hvilken analyselinse en konkret problemstilling tilhører.}
\end{figure}

\subsection{Hovedteori}

\subsubsection*{Strategi som posisjonering (Porter 1996)}
Michael Porter sin klassiske artikkel \emph{What Is Strategy?} er det teoretiske
ankeret. Hovedpoenget er at strategi handler om å \keypoint{utføre andre aktiviteter
enn konkurrentene, eller utføre lignende aktiviteter på en annen måte}. Strategi
betyr å være \emph{annerledes}, ikke bare \emph{bedre}.

\begin{definition}{Strategisk posisjonering}
Et bevisst valg av en unik kombinasjon av aktiviteter som skaper en
verdiblanding ingen konkurrenter kan matche uten å gi opp sin egen posisjon.
\end{definition}

Porter skiller skarpt mellom \term{operasjonell effektivitet} og \term{strategi}:
\begin{itemize}
  \item \textbf{Operasjonell effektivitet} -- å gjøre de samme aktivitetene bedre
    enn konkurrentene. Konvergerer over tid (alle kopierer best practice).
  \item \textbf{Strategi} -- å velge bevisst hvilke aktiviteter man skal gjøre, og
    hvordan disse henger sammen i et unikt system.
\end{itemize}

\subsubsection*{Strategi som valg og avveininger}
Et kjernepoeng: strategi krever \keypoint{trade-offs}. Man kan ikke samtidig være
den billigste og den mest premium aktøren. Et bevisst \emph{fravalg} er like
viktig som et \emph{tilvalg}.

\subsubsection*{Mintzbergs 5 P-er som komplementært perspektiv}
Hvor Porter ser strategi som \emph{posisjon}, gir Mintzberg fem alternative
linser:
\begin{itemize}
  \item \term{Plan} -- en bevisst, intendert handlingsretning.
  \item \term{Ploy} -- et spesifikt grep ment å lure en konkurrent.
  \item \term{Pattern} -- konsistens i adferd over tid (emergent strategi).
  \item \term{Position} -- en bestemt plass i et marked.
  \item \term{Perspective} -- en grunnleggende måte å se verden på.
\end{itemize}
Linsene plasseres på en akse fra \emph{instrumentelle} til \emph{symbolske}
perspektiver. Porter (1996) og Rumelt (2011) ligger på den instrumentelle
siden; den symbolske siden vektlegger identitet, kultur og mimesis.

\begin{figure}[H]
\centering
\includegraphics[width=0.85\linewidth]{lec01-mintzberg-5ps-strategy.png}
\caption{Mintzbergs 5 P-er på et instrumentelt/symbolsk-kontinuum.
Brukes når man skal begrunne hvilken strategilinse man velger som analytisk
ramme.}
\end{figure}

\subsubsection*{Bred vs.\ smal forståelse av digital strategi}
Chen et al.\ (2010) viser at \term{digital strategi} kan forstås \emph{smalt}
(IT-sentrisk, om infrastrukturvalg) eller \emph{bredt} (forretnings-,
organisasjons- og samfunnsnivå --- en gjennomgripende ``digital X'').
Bharadwaj et al.\ (2013) sin definisjon ligger på den brede siden:
digital strategi er en \emph{forretningsstrategi} formulert og utført ved å
utnytte digitale ressurser for å skape differensiell verdi.

\begin{figure}[H]
\centering
\includegraphics[width=0.85\linewidth]{lec01-chen-narrow-vs-broad-digital-strategy.png}
\caption{Chen et al.\ (2010): smal til bred forståelse av digital strategi.
Kurset adopterer den brede forståelsen.}
\end{figure}

\subsubsection*{IT-strategi i kontekst}
IT-strategi henger sammen med tre tilstøtende fag:
\begin{itemize}
  \item \emph{IT-projektledelse} -- styring av leveranser og prosjekter.
  \item \emph{IT-forandringsledelse} -- den menneskelige siden av digital endring.
  \item \emph{Organisationers opbygning og funktion} -- struktur, prosesser og kultur.
\end{itemize}
Skillet mellom økonomi-blok og IT-blok forklarer hvorfor sosiotekniske og
organisatoriske perspektiver er sentrale for å implementere strategi
(jf.\ \autoref{sec:lec07}).

\begin{figure}[H]
\centering
\includegraphics[width=0.75\linewidth]{lec01-it-strategy-in-context.png}
\caption{IT-strategi i krysningen mellom IT-, organisasjons- og
forandringsfagene.}
\end{figure}

\subsubsection*{Antagelser bak en strategi}
Enhver strategi hviler på antagelser om:
\begin{itemize}
  \item \emph{Omgivelsene} -- konkurrenter, kunder, teknologi, regulering.
  \item \emph{Virksomhetens kapabiliteter} -- hva man kan og ikke kan.
  \item \emph{Hva verdi er} -- for hvem, og hvordan den fanges.
\end{itemize}
Strategisk refleksjon innebærer å gjøre disse antagelsene eksplisitte og teste dem.

\subsection{Sentrale fagbegreper}\label{terms:lec01}
\begin{itemize}
  \item \term{Strategy} -- bevisst valg av posisjon og aktivitetssystem.
  \item \term{Operational effectiveness} -- å gjøre eksisterende aktiviteter bedre.
  \item \term{Strategic positioning} -- å være annerledes, ikke bare bedre.
  \item \term{Trade-offs} -- bevisste fravalg som beskytter posisjonen.
  \item \term{Fit} -- hvordan aktiviteter henger sammen og forsterker hverandre.
  \item \term{5 Ps of strategy} -- Mintzbergs fem perspektiver på strategi.
  \item \term{Digital strategy} -- forretningsstrategi som utnytter digitale
    ressurser (Bharadwaj et al.\ 2013).
\end{itemize}

\subsection{Modeller og rammeverk -- bruksveiledning}

\begin{enumerate}
  \item Bruk Porter (1996) når du må argumentere for at noe \emph{er} eller
    \emph{ikke er} strategi (trade-offs, fit, unikt aktivitetssystem).
  \item Bruk Mintzbergs 5 P-er når du må forklare \emph{hvorfor} du velger en
    bestemt analytisk linse (instrumentell vs.\ symbolsk).
  \item Bruk Chen et al.\ + Bharadwaj når du må forsvare at en problemstilling
    er bredere enn ``IT-investering'' -- den er digital strategi.
\end{enumerate}

\subsection{Eksempler og caser}
\begin{eksempelboks}{Southwest Airlines (Porter)}
Porters klassiske eksempel: bevisst valg om kun korte ruter, ingen
bagasje-overføring, ingen måltider. Hele aktivitetssystemet er konsistent med
posisjonen \emph{lavpris, høyfrekvent}. Andre flyselskaper kunne ikke kopiere
enkeltelementer uten å bryte sin egen posisjon -- klassisk \emph{fit}.
\end{eksempelboks}

\begin{eksempelboks}{Coop App (kurseksempel)}
En digital investering må vurderes mot om den forsterker eller undergraver
den eksisterende posisjonen. Coop App skal styrke butikk-først-posisjonen
(omnichannel), ikke konkurrere med den. Sammenheng med
\autoref{sec:lec04} om DT vs.\ ITEOT.
\end{eksempelboks}

\subsection{Eksamensrelevans}
\begin{itemize}
  \item Vær klar til å definere strategi presist (Porter 1996) og skille fra
    operasjonell effektivitet.
  \item Forklar hvorfor trade-offs er nødvendige.
  \item Bruk strategi-som-posisjonering når du argumenterer for at en
    IT-investering er eller ikke er strategisk (jf.\ \autoref{sec:lec04} om
    digital transformasjon vs.\ IT-enabled transformation).
  \item Kunne diskutere antagelser bak en strategi -- dette er en
    metarefleksjon som eksaminator setter pris på.
  \item Kjenne forskjellen på \emph{smal} og \emph{bred} digital strategi
    (Chen et al., Bharadwaj et al.) -- nyttig for å posisjonere oppgavens
    skop.
\end{itemize}
